\documentclass[11pt,a4paper]{article}
% Common preamble for the hint sheets (\input from each file)
\usepackage[utf8]{inputenc}
\usepackage[T1]{fontenc}
\usepackage{lmodern}
\usepackage{amsmath,amssymb,amsthm,mathtools}
\usepackage[a4paper,margin=2.3cm]{geometry}
\usepackage{enumitem}
\usepackage{xcolor}
\usepackage[hidelinks]{hyperref}
\renewcommand{\leq}{\leqslant}
\renewcommand{\le}{\leqslant}
\renewcommand{\geq}{\geqslant}
\renewcommand{\ge}{\geqslant}
\setlength{\parindent}{0pt}
\setlength{\parskip}{0.35em}
% \hintstitle{Sheet title}{Class work / Homework}
\newcommand{\hintstitle}[2]{\begin{center}{\LARGE\bfseries Hints: #1\par}\vspace{0.5em}{\large #2\par}\end{center}
\noindent\textit{The hints to each problem are cumulative. Read one, try again, and open the next only if you are still stuck.}\par\medskip\hrule\medskip}
\newcommand{\problem}[1]{\par\bigskip{\large\bfseries Problem #1}\par\nopagebreak\smallskip}
\newcommand{\subpart}[1]{\par\smallskip\textbf{(#1)}\par\nopagebreak}
\newcommand{\hint}[1]{\par\smallskip\noindent\textcolor{black!70}{\textbf{Hint #1.}}\enspace}
\newcommand{\note}[1]{\par\smallskip\noindent\textit{Note.} #1\par}
\providecommand{\lt}{<}
\providecommand{\gt}{>}
\providecommand{\tightlist}{\setlength{\itemsep}{0pt}\setlength{\parskip}{0pt}}
\emergencystretch=3em

\begin{document}
\hintstitle{Abstract algebra -- 6. The hardest problems}{Class work}

\problem{1}
\hint{1} Do not work with $A$ itself. Choose a non-empty $X\subseteq A$ that minimises the ratio $|X+A|/|X|$, and call the minimum $K$. Then $K\le|A+A|/|A|\le c$, and $|Z+A|\ge K|Z|$ for every non-empty $Z\subseteq A$.
\hint{2} Key lemma (Petridis): for every finite $C\subseteq G$, $|X+A+C|\le K|X+C|$. Iterating with $C=(k-1)A$, $(k-2)A,\dots$ gives $|X+kA|\le K^k|X|$, hence $|kA|\le|X+kA|\le c^k|A|$.
\hint{3} Prove the lemma by induction on $|C|$. Write $C=C'\cup\{\gamma\}$ and $Z=\{x\in X: x+A+\gamma\subseteq X+A+C'\}$. Then $|X+A+C|\le|X+A+C'|+|X+A|-|Z+A|\le K|X+C'|+K|X|-K|Z|$.
\hint{4} On the other side, $X+C$ is the disjoint union of $X+C'$ and $(X\setminus W)+\gamma$, where $W=\{x\in X: x+\gamma\in X+C'\}\subseteq Z$. So $|X+C|\ge|X+C'|+|X|-|Z|$, which completes the induction.

\problem{2}
\hint{1} If the ring is a division ring, it is a finite field $\mathbb F_q$ (Wedderburn). Take $n=q-1$: then all non-zero $n$-th powers equal $1$, and $1+1=1$ is impossible.
\hint{2} If it is not a division ring, there is a non-zero non-unit $a$. Some power $e=a^k$ is idempotent, and $e\ne1$. If $e\ne0$, then $x=e$, $y=1-e$, $z=1$ solve the equation for every $n$.
\hint{3} If $a$ is nilpotent, replace it by a suitable power to get $b\ne0$ with $b^2=0$. For $n\ge2$ take $x=b$, $y=z=1$; for $n=1$ take $x=1$, $y=b$, $z=1+b$.

\problem{3}
\hint{1} Translate the conditions into permutations of $\mathbb Z_n$. Let $\tau(x)=x+1$. Then $f$ is a fixed-point-free involution, and the third condition says $(f\tau)^3=\mathrm{id}$.
\hint{2} Compute signs: $f$ is a product of $n/2$ disjoint transpositions (so $n$ is even), and $\tau$ is an $n$-cycle, of sign $-1$.
\hint{3} $f\tau$ has order dividing $3$, so it is even. Hence $(-1)^{n/2}\cdot(-1)=1$, i.e.\ $n/2$ is odd.

\problem{4}
\hint{1} $\mathbb Q^*$: find a subgroup $H'\le H$ and a monomorphism $H'\to\mathbb Q^*$ that cannot be extended because a root is missing in $\mathbb Q$. Try $H=\mathbb Z\supset H'=2\mathbb Z$ with $\mu'(2k)=2^k$. An extension would need $\mu(1)^2=2$.
\hint{2} $(\mathbb Q,+)$ is divisible. Extend $\mu'$ one element at a time. Given $\mu$ on $L$ and $h\notin L$: if no multiple of $h$ lies in $L$, set $\mu(h)=0$; otherwise let $m>0$ be minimal with $mh\in L$ and set $\mu(h)=\mu(mh)/m$. Check that this is well defined on $L+\mathbb Zh$.
\hint{3} Finish with Zorn's lemma applied to the extensions of $\mu'$ ordered by inclusion (any homomorphism extending $\mu'$ is fine; it need not be injective).

\problem{5}
\hint{1} Show that $H$ always equals $L=\langle$all $2^n$ words of length $n$ in $g,h\rangle$. Clearly $H\subseteq L$, since $H$ is generated by some of these words. Call a length-$n$ word ``good'' if it lies in $H$; every length-$n$ factor of $w$ is good.
\hint{2} Rectangle rule: if $PR$, $SR$, $SQ$ are good (with $|P|=|S|$), then $PQ=(PR)(SR)^{-1}(SQ)$ is good.
\hint{3} Two facts about $w$. Every factor has a right extension. And since $w$ is not eventually periodic, for every length $\ell$ there is a right special factor $C$ of length $\ell$ ($Cg$ and $Ch$ both occur). Otherwise the next letter would be determined by the previous $\ell$ letters, and $w$ would be eventually periodic.
\hint{4} Prove by induction on $r$ that every word $BA$ with $B$ a factor of length $n-r$ and $A$ arbitrary of length $r$ is good. For the step, take $BxA$ with $|A|=r$. Choose $y$ with $By$ a factor and $C$ right special of length $n-r-1$. Then $BxA=(ByA)(CyA)^{-1}(CxA)$. At $r=n$ all words are good.

\end{document}
